\documentclass[12pt]{ctexart}
\usepackage{amsmath}
\usepackage{amssymb}
\usepackage{geometry}
\geometry{a4paper, margin=1in}

\title{铜陵一中2025级高一入学信心卷 \\ 数学}
\author{}
\date{}

\begin{document}

\maketitle

\begin{center}
    满分150分 \quad 时间120分钟 \\
    \textbf{考察范围：}必修一前三章、其他基础初等数学内容
\end{center}

\section*{一、单项选择题（每小题5分，共8小题，40分）}

\begin{enumerate}
    \item 已知集合$A=\{x|2<x<5\}$, 集合$B=\{x|x\ge3\}$, 则$A\cap B=$（ \quad ）
    \begin{enumerate}
        \item[A.] $\{x|3\le x<5\}$
        \item[B.] $\{x|x>2\}$
        \item[C.] $\{x|x\le2\}$
        \item[D.] $\{x|2<x\le3\}$
    \end{enumerate}

    \item 从集合$\{1,2,3,4,5,6,7,8,9\}$中随机地取四个互不相同的数，则其中任意两个数之和均不等于10的概率为（ \quad ）
    \begin{enumerate}
        \item[A.] $\frac{1}{4}$
        \item[B.] $\frac{5}{21}$
        \item[C.] $\frac{8}{21}$
        \item[D.] $\frac{2}{3}$
    \end{enumerate}

    \item 函数$f(x) = \sqrt{2x^2+2} + x$的最小值为（ \quad ）
    \begin{enumerate}
        \item[A.] $2\sqrt{2}$
        \item[B.] $2\sqrt{3}-1$
        \item[C.] $\sqrt{6}$
        \item[D.] $5 + 16$
    \end{enumerate}

    \item 定义在$\mathbb{R}$上的函数$f(x)$满足$f(x+1) = \frac{1}{2} + \sqrt{f(x)-[f(x)]^2}$ ,则$f(0)+f(2017)$的最大可能值为（ \quad ）
    \begin{enumerate}
        \item[A.] $\frac{\sqrt{5}-1}{2}$
        \item[B.] $\frac{\sqrt{6}}{2}$
        \item[C.] $\frac{3\sqrt{3} - 1008}{2017}$
        \item[D.] $1+\frac{\sqrt{2}}{2}$
    \end{enumerate}

    \item 设a、b、c为正数，$a<b$，若a,b为一元二次方程 $ax^2-bx+c =0$的两个根，且a,b,c是一个三角形的三边长，则$a+b-c$的取值范围是（ \quad ）
    \begin{enumerate}
        \item[A.] $(\frac{21}{16}, 2)$
        \item[B.] $(\frac{21}{16}, +\infty)$
        \item[C.] $(\frac{7}{8}, \frac{\sqrt{5}-1}{2})$
        \item[D.] $(\frac{5}{8}, \frac{\sqrt{2}+1}{2})$
    \end{enumerate}

    \item 若函数$f(x)= (2x^5+2x^4-53x^3-57x+54)^{2018}$ （$x\in\mathbb{R}$），则$f\left(\frac{\sqrt{11}-1}{2}\right)=$（ \quad ）
    \begin{enumerate}
        \item[A.] 0
        \item[B.] 1
        \item[C.] 2
        \item[D.] -1
    \end{enumerate}

    \item 已知A、B、C为$\{1,2,3,4,5,6,7\}$的子集，且满足两两交集中元素个数为1，$A\cap B\cap C=\emptyset$，则这样的三元组（A,B,C）（无序）的个数为（ \quad ）
    \begin{enumerate}
        \item[A.] 14400
        \item[B.] 42300
        \item[C.] 53760
        \item[D.] 64230
    \end{enumerate}

    \item 集合$M=\left\{a\in\mathbb{Z} \left| a= \frac{x+y+z}{t} , 3^x+3^y+3^z=3^t ,x,y,z,t\in\mathbb{Z} \right.\right\}$中所有元素之和为（ \quad ）
    \begin{enumerate}
        \item[A.] 6
        \item[B.] 9
        \item[C.] 12
        \item[D.] 15
    \end{enumerate}
\end{enumerate}

\section*{二、多项选择题（每小题6分，共3小题，18分，每个小题全部选对得6分，部分选对得部分分，有错选的得0分）}

\begin{enumerate}
    \setcounter{enumi}{8}
    \item 以下数中，在函数$f(x) = \frac{x(x^2+8)(8-x)}{x+1}$ 的值域内的是（ \quad ）
    \begin{enumerate}
        \item[A.] 0
        \item[B.] 3
        \item[C.] 6
        \item[D.] 9
    \end{enumerate}

    \item 已知数列$\{ a_n \}$满足：$a_1=1， a_2=9， a_{n+2} = 10a_{n+1} - a_n$,其中n为正整数。则下列给出的n中，满足 $a_n$ 不是3的方幂的是（即不存在非负整数k使得 $a_n = 3^k$ ）（ \quad ）
    \begin{enumerate}
        \item[A.] 1
        \item[B.] 18
        \item[C.] 57
        \item[D.] 126
    \end{enumerate}

    \item 如果正整数n使得任意不同于n的正整数m，均有$\{ \sqrt{2}n \} \neq \{ \sqrt{2}m \}$，则称n为“好数”，这里$\{x\}$表示x的小数部分。则下列形式的n中，一定是好数的是（ \quad ）
    （其中，k,t为任意正整数，p为奇素数）
    \begin{enumerate}
        \item[A.] $n= (2^{2k}+1)^2$
        \item[B.] $n= 2^k \cdot 3^t$
        \item[C.] $n= \frac{2^{2k}-1}{3}$
        \item[D.] $n= \frac{3^p+2^p}{5}$
    \end{enumerate}
\end{enumerate}

\section*{三、填空题（每小题5分，共3小题，15分）}

\begin{enumerate}
    \setcounter{enumi}{11}
    \item 若函数$f(x)=2018-ax^2$（$a>0$）的图像与x轴围成的封闭图形内部和边界共有 $2018^2$ 个整点（横纵坐标都是整数的点），则a的取值范围为\underline{\hspace{5cm}}.
    
    \item 已知非负整数a,b,c,d满足 $a^2+b^2+c^2+d^2 =20$,则 $(a+b+c+d)^2 \cdot \left(\frac{1}{a^2+3}+\frac{1}{b^2+3}+\frac{1}{c^2+3}+\frac{1}{d^2+3}\right)$ 的最大值为\underline{\hspace{5cm}}.

    \item 从前2008个正整数构成的集合$M=\{1,2，...，2008\}$中取出一个k元子集A，使得A中任意两个元素之和都不能被它们的差（取正值）整除，则k的最大值为\underline{\hspace{5cm}}.
\end{enumerate}

\section*{四、解答题（本大题共5小题，第15小题13分，第16、17小题15分，第18、19小题17分，共77分，解答应有必要的推理过程和文字说明）}

\begin{enumerate}
    \setcounter{enumi}{14}
    \item (13分) 已知$a>0, 12a+5b+2c>0$，证明：关于x的一元二次方程 $ax^2+bx+c =0$在（2,3）上不可能有两个不同的实数根。

    \item (15分) 设二次函数$f(x)= x^2+bx+c$，已知对任意实数b，均存在实数$x\in[1,2]$,使得不等式$|f(x)|\ge x$成立，求实数c的取值范围。

    \item (15分) 解方程：$\frac{3x}{3+\sqrt{8x-3}} + \frac{3x}{3-\sqrt{8x-3}} = 1$

    \item (17分)
    \begin{enumerate}
        \item[(1)] 已知$x,y>0$，求证：当$x\ge y$时 $\frac{x}{y}\ge\frac{x+1}{y+1}$ ;当$x\le y$时， $\frac{x}{y}\le\frac{x+1}{y+1}$ 。
        \item[(2)] 已知 $a_1,a_2,...,a_n >0$，令 $a_{n+1} = a_1$，证明：
        $$ \frac{a_2}{a_1}+\frac{a_3}{a_2}+...+\frac{a_n}{a_{n-1}}+\frac{a_{n+1}}{a_n} \ge \frac{\sqrt{a_2^2+1}}{\sqrt{a_1^2+1}}+\frac{\sqrt{a_3^2+1}}{\sqrt{a_2^2+1}}+...+\frac{\sqrt{a_n^2+1}}{\sqrt{a_{n-1}^2+1}}+\frac{\sqrt{a_{n+1}^2+1}}{\sqrt{a_n^2+1}} $$
    \end{enumerate}

    \item (17分) 设A、B为正整数，S是由一些正整数构成的一个集合，具有以下性质：
    \begin{enumerate}
        \item[①] 对任意非负整数k，有 $A^k \in S$；
        \item[②] 若正整数$n\in S$，则n的每个正因数均属于S；
        \item[③] 若$m,n\in S$，且m,n互素（m,n的最大公因数为1），则$mn\in S$；
        \item[④] 若$n\in S$，则$An+B\in S$.
    \end{enumerate}
    （已知裴蜀定理：对任意互素的正整数a,b，一定存在整数x,y使得ax+by=1）
    \begin{enumerate}
        \item[(1)] A=1时，证明：任意与B互素的正整数m均属于S；
        \item[(2)] A>1时，证明：若$n\in S$，则 $A^k n + B\frac{A^k-1}{A-1} \in S$ ，其中k为任意正整数；
        \item[(3)] A>1时，证明：与B互素的所有正整数均属于S。
    \end{enumerate}
\end{enumerate}

\end{document}